% !TeX root = ../main.tex
\chapter{Experimental Notebooks}
\label{app:notebooks}

\section{Corpus Preparation and Topic Modeling}
\label{sec:app_nb01}

\begin{lstlisting}[language=Python, caption={Corpus Preparation and Topic Modeling}, label={lst:code_nb01_highlights}]
# Corpus Preparation and Topic Modeling
import pandas as pd
from sklearn.feature_extraction.text import TfidfVectorizer
from sklearn.cluster import MiniBatchKMeans

# Load cleaned statutory and municipal provisions
df = pd.read_json("data/unified_dual_statutory_provisions.jsonl", lines=True)

# Compute c-TF-IDF vector representations across substantive provisions
vectorizer = TfidfVectorizer(max_features=10000, stop_words="english", ngram_range=(1, 2))
X_tfidf = vectorizer.fit_transform(df["provision_text"])

# Cluster provisions into 28 sub-clusters across 8 Macro Legal Domains
kmeans = MiniBatchKMeans(n_clusters=28, random_state=42, batch_size=1024)
df["cluster_id"] = kmeans.fit_predict(X_tfidf)

print(f"Successfully clustered {len(df):,} provisions into {kmeans.n_clusters} clusters.")
\end{lstlisting}

\section{Stage 1 Hybrid Retrieval Benchmarking}
\label{sec:app_nb02}

\begin{lstlisting}[language=Python, caption={Stage 1 Retrieval Sweep and Recall Evaluation}, label={lst:code_nb02_highlights}]
# Stage 1 Hybrid Retrieval Evaluation
import numpy as np

def evaluate_retrieval_recalls(retrieval_engine, benchmark_queries, k_values=[5, 10, 20, 50]):
    """Evaluates candidate retrieval recall across multiple top-k thresholds."""
    recall_results = {k: 0 for k in k_values}
    
    for query in benchmark_queries:
        gold_id = query["gold_statute_id"]
        retrieved = retrieval_engine.retrieve_candidates(query["ordinance_text"], top_k=max(k_values))
        retrieved_ids = [doc["section_id"] for doc in retrieved]
        
        for k in k_values:
            if gold_id in retrieved_ids[:k]:
                recall_results[k] += 1

    total = len(benchmark_queries)
    return {f"Recall@{k}": f"{(hits / total) * 100:.2f}%" for k, hits in recall_results.items()}
\end{lstlisting}

\section{Stage 2 Cross-Encoder Fine-Tuning}
\label{sec:app_nb03}

\begin{lstlisting}[language=Python, caption={Cross-Encoder Fine-Tuning and Validation}, label={lst:code_nb03_highlights}]
# Stage 2 Cross-Encoder Fine-Tuning Loop
from transformers import AutoModelForSequenceClassification, AdamW, get_linear_schedule_with_warmup
from torch.utils.data import DataLoader

model = AutoModelForSequenceClassification.from_pretrained("microsoft/deberta-v3-base", num_labels=3)
optimizer = AdamW(model.parameters(), lr=2e-5, weight_decay=0.01)

# Training loop with validation checkpointing
for epoch in range(5):
    model.train()
    for batch in train_dataloader:
        optimizer.zero_grad()
        outputs = model(**batch)
        loss = outputs.loss
        loss.backward()
        optimizer.step()
        scheduler.step()

    # Validation evaluation for Macro F1 and Conflict F1
    val_macro_f1, val_conflict_f1 = evaluate_nli_fold(model, val_dataloader)
    print(f"Epoch {epoch+1}: Val Macro F1 = {val_macro_f1:.4f}, Conflict F1 = {val_conflict_f1:.4f}")
\end{lstlisting}

\section{Full-Document System Stress Testing}
\label{sec:app_nb04}

\begin{lstlisting}[language=Python, caption={End-to-End System Audit and Verification}, label={lst:code_nb04_highlights}]
# System-Level Full Document Evaluation
def audit_complete_draft_ordinance(draft_pdf_path, pipeline):
    """Executes end-to-end audit: ingestion -> boilerplate filter -> retrieval -> NLI."""
    raw_sections = extract_folio_sections(draft_pdf_path)
    operative_sections = [s for s in raw_sections if not is_boilerplate(s)]
    
    audit_report = []
    for section in operative_sections:
        candidates = pipeline.retrieval.retrieve_candidates(section["text"], top_k=50)
        inference_results = pipeline.inference.classify_candidates(section["text"], candidates)
        audit_report.append({
            "section_num": section["num"],
            "max_conflict_risk": max(r["conflict_risk"] for r in inference_results),
            "flagged": any(r["triage_flag"] for r in inference_results)
        })
    return audit_report
\end{lstlisting}

\section{Archival VLM Batch Digitization Notebook (Google Kaggle)}
\label{sec:app_nb05}

\begin{lstlisting}[language=Python, caption={Headless Kaggle Batch Transcription and Checkpoint Packaging}, label={lst:code_nb05_highlights}]
# Archival VLM Ingestion Loop on Google Kaggle GPU (Tesla T4)
def run_archival_batch_ingestion(input_root, output_dir, output_zip, target_years=None):
    """Executes high-throughput headless VLM batch transcription with fault tolerance."""
    all_pdfs = sorted(p for p in input_root.rglob("*.pdf") if p.is_file())
    if target_years:
        all_pdfs = [p for p in all_pdfs if any(y in p.parts for y in target_years)]

    processed, skipped, total_pages = 0, 0, 0
    for idx, pdf_path in enumerate(all_pdfs, 1):
        out_path = output_dir / pdf_path.relative_to(input_root).with_suffix(".txt")
        if out_path.exists() and out_path.stat().st_size > 100:
            skipped += 1
            continue

        pages = process_pdf_document(pdf_path, out_path)
        processed += 1
        total_pages += pages

        # Periodic checkpoint archive packaging every 25 enactments
        if processed > 0 and processed % 25 == 0:
            create_checkpoint_zip(output_dir, output_zip)

    create_checkpoint_zip(output_dir, output_zip)
    return {"processed": processed, "skipped": skipped, "pages": total_pages}
\end{lstlisting}
